% =======================
% 基本パッケージ
% =======================
\usepackage{float}
\usepackage{titlesec}
\usepackage{setspace}
\usepackage{fontspec}
\usepackage{luatexja}
\usepackage{luatexja-fontspec}
\usepackage{graphicx}
\usepackage{qrcode}
% \usepackage{xcolor}
\usepackage{fancyhdr}
\usepackage{tcolorbox}
\usepackage{xparse}
\usepackage{enumitem}
\usepackage{needspace}
% \usepackage{hyperref}

% tcolorbox ライブラリ
\tcbuselibrary{skins,breakable}

% =======================
% 図の配置
% =======================
\floatplacement{figure}{H}

% =======================
% 行間・段落
% =======================
\onehalfspacing
\setstretch{1.2}
\setlength{\parskip}{0.5em}

% =======================
% フォント設定
% =======================
\setmainfont{Times New Roman}
\setsansfont{Arial}
\setmonofont{Latin Modern Mono}

\setmainjfont{BIZ UDPGothic}
\setsansjfont{BIZ UDPGothic}

\newcommand{\fontsetA}[1]{%
  {\fontspec{TeX Gyre Termes}%
  \jfontspec{BIZ UDPGothic}#1}%
}

\newcommand{\fontsetB}[1]{%
  {\fontspec{BIZ UDPMincho}%
  \jfontspec{BIZ UDPMincho}#1}%
}

\newenvironment{minchoblock}
  {\fontspec{BIZ UDPMincho}\jfontspec{BIZ UDPMincho}}
  {}

\newenvironment{gothicblock}
  {\fontspec{BIZ UDPGothic}\jfontspec{BIZ UDPGothic}}
  {}

% =======================
% 見出し設定
% =======================

\renewcommand{\chaptername}{Chapter}

% \titleformat{\chapter}
%   {\bfseries\fontsize{16pt}{18pt}\selectfont}
%   {\thechapter}
%   {1em}
%   {}

% chapter はゴシック体で表示
% \titleformat{\chapter}
%   {\sffamily\bfseries\fontsize{16pt}{18pt}\selectfont}
%   {\thechapter}{1em}{}

% \titleformat{\chapter}[block]
%   {\bfseries\fontsize{16pt}{18pt}\selectfont}
%   {\thechapter}
%   {1em}
%   {}


%\titlespacing*{\chapter}{0pt}{-20pt}{6pt}
%\titleformat{\chapter}[block]
%  {\gtfamily\sffamily\bfseries\fontsize{18pt}{18pt}\selectfont}
%  {\thechapter}{1em}{}

\titlespacing*{\chapter}{0pt}{70pt}{10pt}     % 章の前後余白 {左}{上}{下}
\titleformat{\chapter}[display]
  {\gtfamily\sffamily\bfseries\filcenter\def\hdbreak{\\}} % 中央寄せ＋見出し組版時のみ {{br}} を強制改行に
  {\fontsize{14pt}{16pt}\selectfont 第\thechapter 章}   % ラベル（独立行・上）
  {1.2em}
  {\fontsize{28pt}{34pt}\selectfont
   \hyphenpenalty=10000\exhyphenpenalty=10000\hbadness=10000\emergencystretch=1em % 単語途中でのハイフン改行を抑止
   \begin{minipage}[b]{0.8\linewidth}\centering}        % タイトルを linewidth の70%幅・中央・折返し可に
  [\vskip 2mm \hrule height 2pt\end{minipage}]          % タイトル直下(2mm)に70%幅の下線（minipage内・行間グルー無し）


\titleformat{\section}[block]
  {\gtfamily\sffamily\bfseries\fontsize{14pt}{16pt}\selectfont\def\hdbreak{\\}} % {{br}} を強制改行に
  {\thesection}{1em}{}

% \titleformat{\section}
%   {\bfseries\fontsize{14pt}{16pt}\selectfont}
%   {\thesection}
%   {1em}
%   {}

\titleformat{\subsection}
  {\bfseries\fontsize{12pt}{14pt}\selectfont\def\hdbreak{\\}} % {{br}} を強制改行に
  {\thesubsection}
  {1em}
  {}

% 章、節、項の前後余白
% \titlespacing{<コマンド>}{<左>}{<前>}{<後>}
% \titlespacing{\chapter}{0pt}{0pt}{12pt}
\titlespacing{\section}{0pt}{18pt}{6pt}
\titlespacing{\subsection}{0pt}{12pt}{4pt}


% \let\oldtextbf\textbf
% \renewcommand{\textbf}[1]{%
%   {\color{blue}\oldtextbf{#1}}%
% }


% =======================
% ヘッダ用マーク設定
% =======================
\makeatletter
\renewcommand{\chaptermark}[1]{%
  \markboth{%
    \ifnum\c@secnumdepth>\m@ne
      \if@mainmatter \@chapapp\ \thechapter.\ \fi
    \fi
    #1%
  }{}%
}
\renewcommand{\sectionmark}[1]{%
  \markright{%
    \ifnum\c@secnumdepth>\z@
      \thesection\quad
    \fi
    #1%
  }%
}
\makeatother

% =======================
% 柱（ヘッダ）設定
% =======================
% 本文ページ（section/subsection を含むページ）の柱に節タイトルを右寄せで出す。
% 上の「ヘッダ用マーク設定」の \sectionmark が \rightmark に節タイトルを入れる。
% chapter の開きページや目次ページは \thispagestyle{plain} なので柱は出ない。
\pagestyle{fancy}
\fancyhead{}
\fancyhead[R]{\gtfamily\sffamily\small\rightmark}   % 右に節タイトル
\renewcommand{\headrulewidth}{0.4pt}
\fancyfoot{}
\fancyfoot[C]{\thepage}                              % ノンブル（copyright 未指定時のフォールバック）
% 開きページ(plain)は柱なし・ノンブルのみ（copyright 指定時はテンプレ側が再定義して上書き）
\fancypagestyle{plain}{%
  \fancyhf{}%
  \renewcommand{\headrulewidth}{0pt}%
  \fancyfoot[C]{\thepage}%
}

% =======================
% 目次設定
% =======================
\renewcommand{\contentsname}{目次}   % 目次見出しの文言（既定 Contents → 目次）

\makeatletter
\renewcommand*\l@section{\@dottedtocline{1}{1.5em}{3.0em}}
\renewcommand*\l@subsection{\@dottedtocline{2}{3.8em}{4.8em}}

% 目次見出しを章扉と別レイアウトにする。
% 注意: number-sections:false（pandoc 既定）では secnumdepth=-\maxdimen となり、通常の
% \chapter も「番号なし」= titlesec の numberless 書式で組まれる。目次見出しも
% \tableofcontents 内部の \chapter* なので、titlesec で分離すると通常章まで巻き込む。
% そのため目次だけは \tableofcontents を再定義し、titlesec の chapter 書式を経由せず
% 独自見出し（中央・小さめ・全幅の細い下線）にする。
\renewcommand\tableofcontents{%
  \if@twocolumn\@restonecoltrue\onecolumn\else\@restonecolfalse\fi
  \clearpage
  \thispagestyle{plain}%
  \@mkboth{\contentsname}{\contentsname}%
  \vspace*{0pt}%
  {\centering\gtfamily\sffamily\bfseries\fontsize{20pt}{24pt}\selectfont\contentsname\par}%
  \vskip 1.5mm \hrule height 0.5pt \vskip 10pt
  \@starttoc{toc}%
  \if@restonecol\twocolumn\fi
}
\makeatother

% =======================
% ページ設定
% =======================
\AtBeginDocument{%
  \setcounter{chapter}{0}
}

% =======================
% 箇条書き
% =======================
\setlist[itemize]{label=\scriptsize●}

% =======================
% 引用ボックス設定
% =======================
% \tcbset{
%   widequote/base/.style={
%     breakable,
%     after upper=\unskip
%   },
%   widequote/indent/.style={
%     widequote/base,
%     blanker,
%     left=4mm,
%     right=2mm
%   },
%   widequote/frame/.style={
%     widequote/base,
%     colback=white,
%     colframe=black,
%     boxrule=0.4pt,
%     left=4mm,
%     right=2mm,
%     top=1.5mm,
%     bottom=1.5mm
%   },
%   widequote/noframe/.style={
%     widequote/base,
%     blanker,
%     left=4mm,
%     right=2mm,
%     top=1.5mm,
%     bottom=1.5mm
%   },
%   quote/frame/.style={
%     breakable,
%     colback=white,
%     colframe=black,
%     boxrule=0.4pt,
%     left=2mm,
%     right=2mm,
%     top=1.5mm,
%     bottom=1.5mm,
%     after upper=\unskip
%   }
% }

% =======================
% インデント環境
% =======================
\newenvironment{blockindent}[1][2em]
{
  \begin{list}{}%
    {
      \setlength{\leftmargin}{#1}%
      \setlength{\rightmargin}{0em}%
      \setlength{\topsep}{0pt}%
      \setlength{\partopsep}{0pt}%
      \setlength{\parsep}{0pt}%
      \setlength{\itemsep}{0pt}%
    }%
  \item[]
}
{
  \end{list}
}

% =======================
% 大きい文字環境
% =======================
\newenvironment{largefontblock}
  {\Large}
  {}
  
  
% =======================
% 引用ボックス・fenced div 用設定
% =======================
\tcbset{
  leadline/font/.style args={#1}{
    fontupper=\gtfamily\sffamily\bfseries #1
  },
  note/font/.style args={#1}{
    fontupper=#1
  },
  warning/font/.style args={#1}{
    fontupper=\gtfamily\sffamily\bfseries #1
  },
  widequote/font/.style args={#1}{
    fontupper=#1
  },
  widequote/stretch/.style args={#1}{
    before upper={\setstretch{#1}}
  },
  leadline/base/.style={
    breakable,
    after upper=\unskip
  },
  leadline/frame/.style={
    leadline/base,
    colback=white,
    colframe=black,
    boxrule=0.4pt,
    left=0mm,
    right=0mm,
    top=0.5mm,
    bottom=0.5mm
  },
  leadline/noframe/.style={
    leadline/base,
    blanker,
    left=0mm,
    right=0mm,
    top=0.5mm,
    bottom=0.5mm
  },
  note/base/.style={
    breakable,
    after upper=\unskip
  },
  note/frame/.style={
    note/base,
    colback=black!3,
    colframe=black!30,
    boxrule=0.3pt,
    left=2mm,
    right=2mm,
    top=1mm,
    bottom=1mm
  },
  note/noframe/.style={
    note/base,
    blanker,
    left=2mm,
    right=2mm,
    top=1mm,
    bottom=1mm
  },
  warning/base/.style={
    breakable,
    after upper=\unskip
  },
  warning/frame/.style={
    warning/base,
    colback=black!5,
    colframe=black,
    boxrule=0.5pt,
    left=2mm,
    right=2mm,
    top=1mm,
    bottom=1mm
  },
  warning/noframe/.style={
    warning/base,
    blanker,
    left=2mm,
    right=2mm,
    top=1mm,
    bottom=1mm
  },
  widequote/base/.style={
    breakable,
    widequote/stretch={1.4},
    after upper=\unskip
  },
  widequote/indent/.style={
    widequote/base,
    blanker,
    left=4mm,
    right=2mm
  },
  widequote/frame/.style={
    widequote/base,
    colback=white,
    colframe=black,
    boxrule=0.4pt,
    left=4mm,
    right=2mm,
    top=1.5mm,
    bottom=1.5mm
  },
  widequote/noframe/.style={
    widequote/base,
    blanker,
    left=4mm,
    right=2mm,
    top=1.5mm,
    bottom=1.5mm
  }
}

% =======================
% quote 環境
% =======================
% \renewenvironment{quote}
%   {
%     \begin{tcolorbox}[quote/frame]
%     \begin{spacing}{1.2}
%   }
%   {
%     \end{spacing}
%     \unskip
%     \end{tcolorbox}
%   }

% =======================
% widequote 環境
% =======================
\NewDocumentEnvironment{widequote}{ O{} }
  {
    \begin{tcolorbox}[widequote/noframe,#1]
    \ignorespaces
    \begin{spacing}{1.4}
  }
  {
    \end{spacing}
    \unskip\ignorespacesafterend
    \end{tcolorbox}
  }

\newenvironment{widequotenoframe}
  {
    \begin{widequote}[widequote/noframe]
  }
  {
    \end{widequote}
  }

% =======================
% fenced div 用ボックス
% =======================
\newtcolorbox{leadlinebox}[1][]{
  enhanced,
  breakable,
  blanker,
  before skip=10pt,
  after skip=8pt,
  left=0pt,
  right=0pt,
  top=1mm,
  bottom=1mm,
  leadline/font={\Large},
  #1
}

\newtcolorbox{notebox}[1][]{
  enhanced,
  breakable,
  colback=black!3,
  colframe=black!35,
  boxrule=0.4pt,
  left=3mm,
  right=3mm,
  top=1.5mm,
  bottom=1.5mm,
  before skip=8pt,
  after skip=8pt,
  note/font={\normalsize},
  #1
}

\newtcolorbox{warningbox}[1][]{
  enhanced,
  breakable,
  colback=black!6,
  colframe=black,
  boxrule=0.8pt,
  left=3mm,
  right=3mm,
  top=1.5mm,
  bottom=1.5mm,
  before skip=10pt,
  after skip=10pt,
  warning/font={\normalsize},
  #1
}

\newtcolorbox{widequotebox}[1][]{
  widequote/font={\normalsize},
  #1
}
